\subsection{已具备的实验条件}

本项目依托湘潭大学商学院经济管理国家级实验教学示范中心（教高函[2013]10号）开展研究，该中心现有固定资产1,600万元，实验用房1,163.28平方米，下辖十二个专业实验室。申请人所在的经济科学方法与智能实验室（Scientific Methods and Intelligence Laboratory for Economics, SMILE）是融合实验经济学核心方法与人工智能技术的前沿研究平台，聚焦于行为科学、信息科学和数字经济领域的系统研究，为本项目提供核心实验支撑。

硬件设施。 经济仿真模拟实训室配备60台计算机终端，可同时容纳60名被试参与实验，满足大规模行为实验的并行运行需求；具备高速稳定的局域网环境，确保oTree等网络实验平台的流畅运行；配有投影仪、音响设备等辅助设施，支持实验指导语的标准化呈现。

软件平台。 已部署oTree最新版本（基于Python），支持本项目所需的多期博弈、角色分配、实时交互等复杂实验设计。同时配备Stata、EViews、SPSS、Python等数据分析软件，以及Python/Mesa仿真建模环境，分别服务于实验数据分析和智能体仿真模块。问卷调查模块已购置Credamo专业版账号，支持大规模在线问卷的发放与回收。数据资源方面，共享使用CNRDS、Wind等数据库，可支持文献计量和准自然实验分析中对国内科研资助数据的检索。

被试资源。 SMILE实验室已通过ORSEE（Online Recruitment System for Economic Experiments）系统建立了包含2,000余名学生的被试数据库，覆盖经济学、管理学、理学、工学等多个学科，可满足本项目实验室实验对多样化被试样本的需求。

\subsection{尚缺少的实验条件和拟解决途径}

条件缺口一：校外被试样本与在线实验平台。 为增强实验结果的外部效度，本项目计划在校内被试的基础上招募部分校外被试（如其他高校学生和在职科研人员），以检验实验结论对非学生群体的稳健性。

拟解决途径：（a）通过Credamo平台招募符合条件的校外在线被试，该平台已在国内实验经济学研究中广泛使用，被试质量可通过注意力检测和一致性检验加以控制。（b）与国内其他高校实验中心（如武汉大学行为与实验研究中心、厦门大学实验经济学实验室）合作开展跨校实验。申请人与上述机构研究者有既有学术联系，已就合作意向进行了初步沟通。

条件缺口二：国际随机资助机制的微观数据。 文献计量和准自然实验分析需要新西兰HRC、瑞士SNSF等机构的项目级资助数据和评审数据。

拟解决途径：（a）公开数据层面，以Web of Science和Scopus的Funding Information字段为基础，结合各机构年报中的资助统计数据，构建准自然实验所需的面板数据集。（b）合作数据层面，通过申请人的国际合作网络（博士后导师Zacharias Maniadis为欧盟Horizon 2020 SInnoPSis项目ERA Chair，与欧洲多国科研政策研究者保持合作）联系数据持有机构，申请受限数据的学术使用权限。（c）替代数据层面，若国际数据获取受限，将转向利用NSFC公开数据库（基金委官网项目查询系统）和CNKI论文数据库，构建国内科研资助效率分析的数据集。

条件缺口三：高性能计算资源。 智能体仿真（ABM）的大规模参数扫描和蒙特卡洛模拟对计算资源有较高要求。

拟解决途径：常规仿真任务可在SMILE实验室现有计算机上完成；大规模参数扫描（涉及数千组参数组合×数百次模拟重复）将利用湘潭大学高性能计算中心资源，或通过云计算平台（如阿里云弹性计算服务）获取按需计算能力，相关费用已纳入项目预算。

\subsection{利用研究基地的计划与落实情况}

本项目主要依托经济管理国家级实验教学示范中心（湘潭大学）开展实验室实验和数据分析工作。该中心是教育部认定的国家级教学示范中心（教高函[2013]10号），具备完善的实验教学和科研支撑体系。SMILE实验室作为其下属的前沿研究平台，在实验经济学方法与行为科学交叉研究方面具有明确的定位和稳定的运行保障。

在校外合作基地方面，申请人计划通过以下渠道拓展研究资源：（a）塞浦路斯大学SInnoPSis中心——申请人曾在该中心完成博士后研究（2022—2025），与中心主任Zacharias Maniadis教授保持持续合作关系。该中心不仅可提供接入欧盟科研政策研究网络和国际数据资源的渠道，更重要的是，申请人计划与该校合作，在其校内科研经费分配的部分院系评审中试点部分随机机制，构建校级层面的DID自然实验（详见"风险一"应对措施），这将是本项目田野实验的重要实施路径之一。（b）学术会议评审合作——申请人计划联系相关学术会议组织者，参照Goldberg等（2025）利用NeurIPS和ICLR真实评审数据的思路，探索在部分会议的论文接受环节试行随机机制并收集评审数据。（c）国内合作机构——申请人曾在微众银行金融科技学院（深圳南特商学院）担任副研究员（2020—2022），专职从事实验经济学与行为经济学研究，积累了与金融科技机构合作开展行为实验的经验；此外，申请人与武汉大学和厦门大学等高校的实验经济学研究团队保持学术联系，可在被试招募、跨校实验和数据共享方面开展合作。